\documentclass{situationist-preprint}

\title{Luck and Desert in Meritocracy}
\msid{POL-19-a}
\msfamily{POL-19}
\msdiscipline{Political Philosophy}
\mstype{Principal manuscript}
\mskeywords{meritocracy; desert; luck; natural lottery; effort; contribution; responsibility; difference principle}

\begin{document}
\maketitle

\begin{abstract}
Meritocracy holds that social and economic advantages should go to people in proportion to their merit, and that a distribution so made is just because the meritorious deserve what they receive. We argue that meritocracy is not just. Merit, the combination of talent and effort that issues in valued performance, is partly a product of luck: of natural endowment, of the circumstances in which capacities are developed, and of the market's demand for the capacities one happens to have. Luck is not deserved. We state a parity principle for desert, on which persons who agree in everything under their control deserve the same rewards, and prove that any meritocratic distribution violates it whenever merit varies with luck, which it always does. The principle is the distributive counterpart of the control principle for culpability defended in ETH-13, and like that principle it leaves responsibility intact: the talented are responsible for their achievements, and achievements are theirs, but the size of a reward is not settled by the scope of responsibility. The argument does not condemn selecting people for positions by their qualifications; it denies that such selection is justified by desert. Its justification, where it has one, lies in the principles of POL-01.
\end{abstract}

\section{The problem}

The word \enquote{meritocracy} was coined by \citet{Young1958} for a dystopia, a society in which intelligence plus effort determines position and the successful believe they have earned everything they have. It has since become the name of an ideal. On the meritocratic view, careers should be open to talents, positions should go to the most qualified, and income and esteem should follow performance; a society in which they do is just, because each person receives what she deserves.

\citet[\S\S 17, 48]{Rawls1971} rejected the desert claim. The distribution of natural talents is, in his phrase, arbitrary from a moral point of view, and even the willingness to make an effort depends in part on fortunate family and social circumstances; no one deserves his greater natural capacity or his more favourable starting place. Rawls replaced desert, at the level of the basic structure, with legitimate expectations: people are entitled to what just institutions lead them to expect, but the institutions are not designed to reward desert. \citet[ch.~7]{Nozick1974} replied that the foundations of desert need not themselves be deserved all the way down, and that people are entitled to their natural assets whether or not they deserve them. \citet{Sher1987} and \citet{Miller1999} developed theories of desert on which the hardworking and those who contribute more deserve more, and \citet{Feinberg1970} gave the classic analysis of desert as a relation between a person, a mode of treatment and a basis in some fact about the person.

Recent work has returned to the question with empirical as well as normative arguments. \citet{Frank2016} documents the role of chance in economic success, especially in winner-take-all markets where small differences in performance, or in timing, produce very large differences in reward. \citet{Sandel2020} argues that the meritocratic ideal corrodes the self-understanding of both winners and losers, and \citet{Markovits2019} that meritocratic competition reproduces an elite.

We argue for the conclusion of Rawls's argument by a route that does not require his premise that nothing is deserved. Merit is partly a product of luck, and luck is not deserved; meritocracy rewards merit as a whole, and so rewards luck as if it were deserved.

\section{Framework}

\begin{definition}\label{d:merit}
The \emph{merit} of a person, relative to a social role or market, is the value of her performance in it, as it results from her talents, the development of those talents, and her effort.
\end{definition}

\begin{definition}\label{d:meritocracy}
\emph{Meritocracy} is the view that social and economic advantages should be distributed in proportion to merit, and that a distribution so made is just because each person deserves the share that her merit earns.
\end{definition}

Definition~\ref{d:meritocracy} picks out meritocracy as a theory of distributive justice. It is to be distinguished from \emph{merit-based selection}, the practice of filling positions with those best qualified to perform them. Selection can be justified on grounds that make no appeal to desert, such as the interest of the people the position serves; meritocracy claims that the rewards attached to positions are deserved by those who occupy them.

\begin{definition}\label{d:luck}
A factor is a matter of \emph{luck} for a person if it is not under her control. Control is control over one's acts in the sense of ETH-13 and ACT-01: an act is controlled when it issues from a reasons-responsive mechanism that is the agent's own.
\end{definition}

\begin{assumption}[Merit is partly luck]\label{a:luck}
Merit depends on factors that are matters of luck: natural endowment, the circumstances of a person's upbringing and education, the demand for her capacities at the time and place she offers them, and the performances of others with whom she competes or cooperates. In every actual society there are persons who agree in all factors under their control and differ in merit.
\end{assumption}

Assumption~\ref{a:luck} is empirical, and its first part is not seriously disputed: nobody chooses his genes or his parents, and nobody controls the price that others will pay for what he can do \citep{Frank2016}. The second part follows if the factors of luck vary independently of the controlled ones, which they do: two equally conscientious students, one of whom had a well-stocked library at home and the other none, will not in general perform equally.

\begin{principle}[Luck is not deserved]\label{p:notdeserved}
No one deserves the factors that are matters of luck for her, and no one deserves a greater reward than another in virtue of a difference in such factors alone.
\end{principle}

\begin{principle}[Parity]\label{p:parity}
If two persons agree in all factors under their control that bear on a reward, they deserve the same reward.
\end{principle}

Principle~\ref{p:parity} follows from Principle~\ref{p:notdeserved}: if two persons agree in everything under their control, any difference between them is a difference in luck, and by Principle~\ref{p:notdeserved} no difference in desert can rest on it. It is the distributive counterpart of the principle that ETH-13 calls the control of culpability: that an agent's culpability depends only on factors under her control together with the evidence available to her. The degree to which a person is blameworthy and the degree to which she deserves a reward are both appraisals of what she has made of what she controlled.

Principle~\ref{p:parity} does not require that desert bases be deserved. It permits a person to deserve a reward for her effort, although she did not deserve the capacity to make it; it requires only that persons who are alike in what they controlled be alike in what they deserve.

\section{The argument}

\begin{lemma}\label{l:differ}
In every actual society, a meritocratic distribution gives different rewards to some persons who agree in all factors under their control.
\end{lemma}

\begin{proof}
By Assumption~\ref{a:luck}, there are persons who agree in all factors under their control and differ in merit. By Definition~\ref{d:meritocracy}, a meritocratic distribution rewards them in proportion to their merit, and so differently.
\end{proof}

\begin{theorem}\label{t:unjust}
Meritocracy is not just: in every actual society, a meritocratic distribution is not a distribution according to desert, and the meritocratic claim that it is just because deserved is false.
\end{theorem}

\begin{proof}
By Lemma~\ref{l:differ}, a meritocratic distribution gives different rewards to some persons who agree in all factors under their control. By Principle~\ref{p:parity}, those persons deserve the same reward. So the distribution gives at least one of them a reward that she does not deserve, or withholds from one a reward that she does. Meritocracy, by Definition~\ref{d:meritocracy}, claims that the distribution is just because each receives what she deserves; that claim is false.
\end{proof}

Theorem~\ref{t:unjust} refutes meritocracy on its own terms. A second argument reaches the same conclusion from the theory of justice defended in POL-01. On that theory, the justice of the basic structure depends on whether it satisfies a principle of equal basic liberty and the modified difference principle, which ranks distributions by the level of the worst-off position and then by the total. A meritocratic distribution is determined by the distribution of merit, and there is no reason why it should maximize the position of the worst off: those with little merit, through bad luck in endowment or circumstance, receive little, whatever alternative arrangements would give them.

\begin{corollary}\label{c:select}
Merit-based selection is not shown to be unjust; its justification, where it has one, does not rest on desert.
\end{corollary}

\begin{proof}
Theorem~\ref{t:unjust} concerns the claim that rewards are deserved. Selection by qualification is justified, where it is, by the interests of those whom the position serves, and inequalities in the rewards attached to positions are justified, on the view of POL-01, if they are to the benefit of the worst off or, at equal benefit to them, increase the total. Neither justification appeals to desert.
\end{proof}

\section{Luck, responsibility and desert}

Our argument may seem to prove too much. If luck undermines the claim of the talented to their rewards, it seems to undermine the claim that anyone is responsible for anything, since every act depends on factors beyond the agent's control. ETH-13 argues that moral luck is real and does not undermine responsibility, and we accept its argument; we must show that ours is consistent with it.

ETH-13 distinguishes two dimensions of responsibility. The \emph{scope} of an agent's responsibility, the set of things she is responsible for, extends to the outcomes of her controlled acts, and depends on luck. Her \emph{culpability}, the degree to which she is blameworthy, depends only on what she controlled. The analogous distinction applies to the successful. A talented and fortunate surgeon is responsible for the lives she saves: they are within the scope of her responsibility, she can rightly take pride in them, and her patients owe her gratitude for what she did. Constitutive and circumstantial luck do not remove her control over her acts, and so they do not remove her responsibility or her authorship of her achievements. What luck does affect is the comparative evaluation of her as an agent: she is no more admirable, as an exercise of will, than an equally conscientious surgeon who, through bad luck in training or in the cases that came to her, saved fewer.

Distributive rewards answer to the second question, not the first. A reward that is claimed as deserved is claimed as a fitting response to the person's worth as a contributor, and on the meritocratic view it varies in magnitude with that worth, as the intensity of blame varies with culpability. That is why Principle~\ref{p:parity} governs it. The scope of responsibility governs what one is answerable for and what one can be credited with; it does not fix how large a share of the social product one deserves. There is thus no conflict between Theorem~\ref{t:unjust} and the thesis of ETH-13: luck leaves the talented responsible for their achievements, and it leaves their achievements undeserved as grounds for unequal reward.

Principle~\ref{p:notdeserved} also underwrites the claim of ETH-22 that the endowment fixed by the natural lottery is not deserved.

\section{Objections}

\begin{objection}[Desert bases need not be deserved]
\citet[ch.~7]{Nozick1974} objected to Rawls that a person can deserve something on the basis of a feature she does not deserve; if desert required deserved bases, nothing would be deserved, since the regress must end in something unearned.
\end{objection}

\begin{reply}
We accept the point and do not rely on the premise it refutes. Principle~\ref{p:parity} allows undeserved bases, such as the capacity for effort, to ground desert. What it forbids is a difference in desert between persons who are alike in what they controlled. The regress is avoided because the principle compares persons; it does not ask whether the base of each person's desert was itself deserved.
\end{reply}

\begin{objection}[Effort]
Meritocracy can be purified: reward effort, which is under control, rather than talent.
\end{objection}

\begin{reply}
A distribution in proportion to effort as exercised under control would satisfy Principle~\ref{p:parity}, but it would not be meritocracy as defined, since merit is the value of performance, not the effort that went into it. It would also be unworkable: effort cannot be observed apart from performance, and the capacity to make it is itself shaped by luck \citep[\S 17]{Rawls1971}. Institutions that reward performance reward luck along with effort, and no institution can separate them.
\end{reply}

\begin{objection}[Contribution]
People deserve in proportion to their contribution to the social product: those who contribute more have done more for others, and a reward that matches contribution is a fitting return \citep{Miller1999}.
\end{objection}

\begin{reply}
The value of a contribution depends on demand, on the contributions of others with which it is combined, and on the number of others able to make it, all of which are matters of luck for the contributor. Two people who do the same work with the same care can contribute very different amounts because one of them works where the work is scarce. Principle~\ref{p:parity} implies that the difference in contribution, so far as it is due to these factors, grounds no difference in desert. The fitting return to a contribution, as a response to what the person did, is gratitude and credit, which follow the scope of her responsibility.
\end{reply}

\begin{objection}[Incentives]
Without rewards for merit, people would not develop or use their talents, and everyone, including the worst off, would be worse off.
\end{objection}

\begin{reply}
The objection supports Corollary~\ref{c:select}. Unequal rewards may be justified as incentives, and the modified difference principle of POL-01 permits them when they benefit the worst off. Such rewards are legitimate expectations under just institutions, not deserts; they are justified by what they do for others, and they are limited by the same consideration.
\end{reply}

\section{What the result does not establish}

The argument does not show that no one deserves anything, that praise and gratitude are never fitting, or that positions should not be filled by the most qualified. It does not show that luck undermines responsibility, which ETH-13 denies. It shows that the rewards of a meritocratic distribution are not deserved and that the justice of a distribution cannot rest on the desert of merit. The companion manuscript POL-19-b turns to merit in practice.

\section{Conclusion}

Merit is the value of a person's performance, and it depends on her endowment, her upbringing and the demand for what she can do, none of which she controls. Persons alike in everything they control deserve alike; a meritocratic distribution treats them differently; so it does not distribute according to desert, and its claim to justice fails. Luck leaves the talented responsible for their achievements, but it does not make their rewards deserved. Merit-based selection and unequal rewards can be justified, where they are, by the principles of POL-01, not by desert.

\begin{thebibliography}{9}\small
\bibitem[Feinberg(1970)]{Feinberg1970} Feinberg, J. (1970). Justice and personal desert. In \emph{Doing and Deserving: Essays in the Theory of Responsibility}. Princeton University Press.
\bibitem[Frank(2016)]{Frank2016} Frank, R.~H. (2016). \emph{Success and Luck: Good Fortune and the Myth of Meritocracy}. Princeton University Press.
\bibitem[Markovits(2019)]{Markovits2019} Markovits, D. (2019). \emph{The Meritocracy Trap}. New York: Penguin Press.
\bibitem[Miller(1999)]{Miller1999} Miller, D. (1999). \emph{Principles of Social Justice}. Harvard University Press.
\bibitem[Nozick(1974)]{Nozick1974} Nozick, R. (1974). \emph{Anarchy, State, and Utopia}. New York: Basic Books.
\bibitem[Rawls(1971)]{Rawls1971} Rawls, J. (1971). \emph{A Theory of Justice}. Harvard University Press.
\bibitem[Sandel(2020)]{Sandel2020} Sandel, M.~J. (2020). \emph{The Tyranny of Merit: What's Become of the Common Good?} New York: Farrar, Straus and Giroux.
\bibitem[Sher(1987)]{Sher1987} Sher, G. (1987). \emph{Desert}. Princeton University Press.
\bibitem[Young(1958)]{Young1958} Young, M. (1958). \emph{The Rise of the Meritocracy, 1870--2033}. London: Thames and Hudson.
\end{thebibliography}

\end{document}
